\documentclass{article}
\usepackage[utf8]{inputenc}
\usepackage{amsmath,amssymb,amsthm}
\usepackage{algorithm,algorithmic}
\usepackage{graphicx}
\usepackage{hyperref}
\usepackage{booktabs}
\usepackage{natbib}
\usepackage[margin=1in]{geometry}
\usepackage{xcolor}

\DeclareMathOperator{\argmin}{arg\,min}

\newtheorem{theorem}{Theorem}
\newtheorem{proposition}{Proposition}
\newtheorem{lemma}{Lemma}
\newtheorem{remark}{Remark}

\title{Block-Wise Differentiable Sinkhorn Attention:\\Scaling Optimal Transport to Massive Contexts}

\author{
Dylan B. Forde \\
\texttt{forde.dylan@gmail.com}
}

\date{}

\begin{document}

\maketitle

\begin{abstract}
We present a memory-efficient, fully differentiable Sinkhorn attention mechanism that scales entropic optimal transport (OT) to massive contexts on TPU hardware. While optimal transport provides a principled, theoretically-grounded framework for discovering structure in sequence data, its application has been severely bottlenecked by quadratic memory requirements and divergent gradient approximations during backpropagation. Our key contribution is a novel \textit{multi-transport-plan backward pass} that correctly differentiates through truncated Sinkhorn iterations by using distinct transport matrices at each chain-rule step. This resolves the fundamental limitation of prior implicit differentiation approaches whose gradients fail to converge in practice. The forward pass uses block-wise Pallas kernels operating in $O(L)$ memory, while the backward pass computes exact gradients \textbf{through the truncated unrolling} using four distinct transport plans, all without materializing the $L \times L$ score matrix. We demonstrate our approach's flexibility on the homologous protein family alignment benchmark from Pfam, achieving 72\% alignment accuracy at 4K tokens vs.\ 20\% for standard softmax attention under identical conditions; a $3.6\times$ improvement that highlights the critical advantage of \textbf{unbalanced, mass-preserving transport normalization}. At 16K+ tokens with banded computation ($W\!=\!2048$), our method trains stably to completion while standard attention exceeds available memory. Ultimately, our work bridges optimal transport theory and hardware-aware kernel design to enable differentiable alignment at scales previously inaccessible to OT methods.
\end{abstract}

\section{Introduction}

Pairwise alignment and structural matching are cornerstones of machine learning applied to sequence data. Classical algorithms such as Smith-Waterman \citep{smith1981identification} provide optimal local alignments but are non-differentiable, precluding their use as trainable components within end-to-end neural networks.

Recent work has explored differentiable relaxations of alignment, notably optimal transport formulations \citep{peyre2019computational}. The Sinkhorn algorithm \citep{cuturi2013sinkhorn} provides a natural bridge: it computes an approximate mass-preserving transport plan via iterative row-column normalization, yielding an output that serves as a robust, differentiable ``soft alignment'' matrix with strong theoretical guarantees. For sequence modeling where gaps are common (such as biological sequences or noisy time-series), we utilize an \textit{unbalanced} Sinkhorn formulation that eschews strict doubly-stochastic properties to allow unaligned mass to be absorbed, preserving local structural relationships without forcing total global bijection.

However, two fundamental challenges have limited the widespread adoption of Optimal Transport and Sinkhorn-based attention in practice:

\begin{enumerate}
    \item \textbf{Memory scaling.} The score matrix $S \in \mathbb{R}^{L \times L}$ required by standard Sinkhorn implementations grows quadratically with sequence length. At $L = 16{,}384$, this single matrix requires 1\,GB in float32; and the backward pass necessitates storing additional intermediates, which quickly exhausts accelerator memory.
    
    \item \textbf{Backward pass correctness.} Differentiating through Sinkhorn iterations is highly non-trivial. Implicit differentiation via the adjoint system \citep{luise2018differential} requires solving a near-singular linear system that often diverges in practice. Conversely, full unrolled backpropagation requires storing all intermediate iterates, thereby exacerbating the memory bottleneck.
\end{enumerate}

We address both challenges simultaneously. Our forward pass utilizes block-wise TPU kernels written in JAX Pallas to compute Sinkhorn potentials $u, v$ via tiled logsumexp reductions, never materializing the full score matrix and maintaining an $O(L)$ memory footprint. For the backward pass, we introduce a novel \textbf{multi-transport-plan truncated backpropagation}. This approach runs two additional differentiable Sinkhorn iterations and correctly propagates gradients using four distinct transport matrices at each chain-rule step; all computed block-wise without exceeding $O(L)$ memory per block pair.

We validate our approach on the Pfam protein family alignment benchmark, demonstrating:
\begin{itemize}
    \item 72\% alignment accuracy at 4K tokens, a $3.6\times$ improvement over the 20\% achieved by softmax attention under identical conditions.
    \item Successful training at 16K tokens with banded computation ($W\!=\!2048$); a regime where standard softmax attention exceeds available memory (OOM).
    \item Stable training without the gradient explosions characteristic of implicit differentiation approaches.
\end{itemize}

\section{Background}

\subsection{Entropic Optimal Transport and Sinkhorn Iterations}

Given two discrete distributions $\mu, \nu$ over $m$ and $n$ points respectively, and a cost matrix $C \in \mathbb{R}^{m \times n}$, the entropically regularized optimal transport problem seeks:
\begin{equation}
    P^* = \argmin_{P \in \Pi(\mu, \nu)} \langle C, P \rangle - \varepsilon H(P)
\end{equation}
where $\Pi(\mu, \nu)$ is the set of transport plans with marginals $\mu, \nu$, and $H(P) = -\sum_{ij} P_{ij} \log P_{ij}$ is the entropy.

The Sinkhorn algorithm solves this via alternating projections in the log domain. Defining $S = -C/\varepsilon$ as the similarity matrix, the iterations maintain dual potentials:
\begin{align}
    u_i &\leftarrow -\text{logsumexp}_j(S_{ij} + v_j) \label{eq:u_update} \\
    v_j &\leftarrow -\text{logsumexp}_i(S_{ij} + u_i) \label{eq:v_update}
\end{align}
The transport plan is subsequently recovered as $P_{ij} = \exp(S_{ij} + u_i + v_j)$.

For biological sequence alignment, the transport plan must accommodate insertions and deletions (gaps). We adopt the unbalanced formulation with a gap score parameter $g = 0$:
\begin{align}
    u_i &\leftarrow -\log\!\Big(\sum_j \exp(S_{ij} + v_j) + \exp(g)\Big) \\
    v_j &\leftarrow -\log\!\Big(\sum_i \exp(S_{ij} + u_i) + \exp(g)\Big)
\end{align}
The gap term allows probability mass to be ``absorbed,'' effectively modeling unaligned positions. We set $g=0$ (meaning $\exp(g)=1$) such that gap absorption is placed on equal footing with an average naive zero-centered alignment score $S_{ij}$, providing a neutrally-biased sink for unaligned mass.

\subsection{Sinkhorn Attention}

In our architecture, we parameterize $S = QK^\top \cdot \alpha / \sqrt{d}$ where $Q, K \in \mathbb{R}^{B \times H \times L \times d}$ are query and key projections, $\alpha = 3.0$ is a sharpness parameter, and $d$ is the head dimension. The output is formulated as:
\begin{equation}
    O = P \cdot V = \exp(S + u \mathbf{1}^\top + \mathbf{1} v^\top) \cdot V
\end{equation}
This mechanism replaces standard softmax attention $O = \text{softmax}(QK^\top / \sqrt{d}) \cdot V$ with a transport-plan-weighted combination.

\section{Method}

\subsection{Block-Wise Forward Pass}

The forward pass computes the Sinkhorn potentials $u, v$ and the output $O = PV$ entirely in a block-wise fashion, conceptually parallel to the tiling strategy of FlashAttention \citep{dao2022flashattention}.

\textbf{Computing $u, v$}: For each query block $Q_i$ of size $B_q$, we scan over all key blocks $K_j$ to compute tiled logsumexp reductions:
\begin{equation}
    u_i = -\log\!\Big(\sum_{j} \exp(Q_i K_j^\top \cdot \alpha/\sqrt{d} + v_j) + \exp(g)\Big)
\end{equation}
Each $S_{ij}$ block is computed on-the-fly and discarded immediately after the reduction, requiring only $O(B_q \times B_k)$ SRAM per block pair.

\textbf{Computing Output}: Similarly, $O_i = \sum_j P_{ij} V_j$ is accumulated block-wise, where $P_{ij} = \exp(S_{ij} + u_i + v_j)$ is recomputed per block.

This strategy yields an $O(L)$ HBM footprint while maintaining $O(L^2 / B_q)$ FLOPs.

\subsection{Why Implicit Differentiation Fails}
\label{sec:implicit_fails}

At the Sinkhorn fixed point, the potentials satisfy:
\begin{equation}
    u = f(S, v), \qquad v = g(S, u)
\end{equation}
Implicit differentiation of this fixed-point system yields the adjoint equations:
\begin{align}
    \lambda_u &= g_u - P \lambda_v \\
    \lambda_v &= g_v - P^\top \lambda_u
\end{align}
where $g_u = \sum_j P_{ij} Z_{ij}$, $g_v = \sum_i P_{ij} Z_{ij}$, and $Z = \frac{\partial L}{\partial O} \cdot V^\top$.

While this linear system can be solved iteratively (e.g., via Gauss-Seidel), the spectral radius of the iteration matrix is $\rho(P P^\top) \approx 1$ for near-doubly-stochastic plans. In practice, this causes:
\begin{itemize}
    \item \textbf{Divergence}: With many iterations, the Gauss-Seidel solver oscillates, producing explosive gradients (manifesting as sudden loss spikes from $\sim 2.8$ to $>4.0$).
    \item \textbf{Ineffective Damping}: Adding $\alpha$-damping ($\lambda \leftarrow \alpha \lambda_{\text{new}} + (1-\alpha)\lambda_{\text{old}}$) prevents divergence but yields a solution too inaccurate for meaningful training.
    \item \textbf{Problem Alteration via Truncation}: Using only 2 Gauss-Seidel steps provides a biased approximation of the \textit{implicit} gradient, which is fundamentally different from the true \textit{truncated backpropagation} gradient.
\end{itemize}

Empirically, all implicit differentiation variants (20-step, damped, 2-step) failed to exceed 14\% alignment accuracy in our testing, whereas our method achieves 72\%.

\subsection{Multi-Transport-Plan Truncated Backpropagation}
\label{sec:multi_P}

Our core insight is that \textbf{truncated backpropagation through $k$ Sinkhorn iterations} provides correct, stable gradients; provided that each chain-rule step utilizes the \textbf{appropriate transport matrix}.

\textbf{Forward (for backward)}: Starting from stopped-gradient potentials $u^{(0)}, v^{(0)}$ (approximating the fixed point after $n$ forward iterations), we recompute $k=2$ additional iterations:
\begin{align}
    u^{(1)} &= f(S, v^{(0)}), & v^{(1)} &= g(S, u^{(1)}) \\
    u^{(2)} &= f(S, v^{(1)}), & v^{(2)} &= g(S, u^{(2)})
\end{align}

\textbf{Backward}: The total gradient $\partial L / \partial S$ receives contributions from five distinct sources along the computational graph. Crucially, although there are five gradient paths, the first two paths both consume the output of the final refinement iteration ($v^{(2)}$) and thus correctly share the identical transport matrix $P^{(2,2)}$. The remaining steps descend the unrolled chain, accumulating gradients utilizing four distinct transport plans overall:

\begin{enumerate}
    \item \textbf{Output}: $\Delta S^{(A)}_{mn} = P^{(2,2)}_{mn} \cdot Z_{mn}$ where $P^{(a,b)} = \exp(S + u^{(a)} + v^{(b)})$
    \item \textbf{$v^{(2)}$ backward}: $\Delta S^{(B)}_{mn} = -g_{v,n} \cdot P^{(2,2)}_{mn}$ \quad {\small (Shares $P^{(2,2)}$ with output!)}
    \item \textbf{$u^{(2)}$ backward}: $\Delta S^{(C)}_{mn} = -\delta u^{(2)}_m \cdot P^{(\mathbf{2,1})}_{mn}$ \quad {\small (Uses $v^{(1)}$, not $v^{(2)}$!)}
    \item \textbf{$v^{(1)}$ backward}: $\Delta S^{(D)}_{mn} = -\delta v^{(1)}_n \cdot P^{(\mathbf{1,1})}_{mn}$ \quad {\small (Uses $u^{(1)}$, not $u^{(2)}$!)}
    \item \textbf{$u^{(1)}$ backward}: $\Delta S^{(E)}_{mn} = -\delta u^{(1)}_m \cdot P^{(\mathbf{1,0})}_{mn}$ \quad {\small (Uses $v^{(0)}$!)}
\end{enumerate}

where the intermediate gradient vectors are defined as:
\begin{align}
    \delta u^{(2)} &= g_u - P^{(2,2)} g_v & &\text{(backward through $v^{(2)}$)} \\
    \delta v^{(1)} &= -[P^{(2,1)}]^\top \delta u^{(2)} & &\text{(backward through $u^{(2)}$)} \\
    \delta u^{(1)} &= -P^{(1,1)} \delta v^{(1)} & &\text{(backward through $v^{(1)}$)}
\end{align}

\begin{proposition}\label{prop:distinct_P}
The four transport matrices $P^{(2,2)}, P^{(2,1)}, P^{(1,1)}, P^{(1,0)}$ are strictly distinct whenever the Sinkhorn potentials have not perfectly converged (i.e., $u^{(0)} \neq u^{(1)} \neq u^{(2)}$ or $v^{(0)} \neq v^{(1)} \neq v^{(2)}$). Replacing all four with a single $P$ produces biased gradients that prevent training convergence. (Proof in Appendix~\ref{app:proof}.)
\end{proposition}

This is confirmed empirically: using a single $P$ for all steps yields $\leq 14$\% accuracy, while utilizing the correct four $P$ matrices yields $72$\%.

\subsection{Block-Wise Backward Implementation}

All backward computations are executed block-wise to maintain the $O(L)$ memory footprint:

\textbf{Phase 1} (Recompute): Run 2 block-wise iterations from $(u^{(0)}, v^{(0)})$ to obtain $u^{(1)}, v^{(1)}, u^{(2)}, v^{(2)}$. Each iteration uses the same tiled scan as the forward pass.

\textbf{Phase 2} (Intermediate Gradients): Compute $g_u, g_v, \delta u^{(2)}, \delta v^{(1)}, \delta u^{(1)}$ via block-wise matrix-vector products. Each matvec uses the correct transport matrix (parameterized by the appropriate $u, v$ pair).

\textbf{Phase 3} (Accumulate $dQ, dK, dV$): For each query-block/key-block pair $(i, j)$:
\begin{enumerate}
    \item Recompute $S_{ij} = Q_i K_j^\top \cdot \alpha/\sqrt{d}$
    \item Compute all four $P^{(a,b)}_{ij}$ using the block-local $u, v$ slices
    \item Accumulate $\Delta S_{ij} = \sum_{k \in \{A..E\}} \Delta S^{(k)}_{ij}$
    \item $dQ_i \mathrel{+}= \Delta S_{ij} \cdot K_j \cdot \alpha/\sqrt{d}$, \quad $dK_j \mathrel{+}= \Delta S_{ij}^\top \cdot Q_i \cdot \alpha/\sqrt{d}$
    \item $dV_j \mathrel{+}= [P^{(2,2)}_{ij}]^\top \cdot dO_i$
\end{enumerate}

Memory per block pair: $5 \times O(B_q \times B_k)$ for the transport matrices, plus $O(B_q \times d)$ and $O(B_k \times d)$ for accumulation. Total HBM footprint: $O(L)$.

\section{Architecture}

We employ a Dense ResNet Aligner architecture for protein sequence alignment:

\textbf{Embedding}: Amino acid tokens (20 standard + gap + pad) are embedded into $d_{\text{model}} = 64$ dimensions using learned positional encodings.

\textbf{Feature Extraction}: Three residual blocks comprising dense layers, LayerNorm, and GELU activations produce contextualized representations.

\textbf{Sinkhorn Attention}: Multi-head attention ($H = 4$, $d_k = 16$) with our Sinkhorn kernel replaces standard softmax attention. The resulting transport plan serves as a soft alignment between ancestor and descendant sequences.

\textbf{Alignment Prediction}: The aligned representations are projected back to token space via a linear head and trained with cross-entropy loss against ground-truth alignments curated from Pfam.

\section{Experiments}

\subsection{Setup}

\textbf{Data}: We benchmark on 1{,}000 homologous protein families from the Pfam database. Ground truth alignments are computed by first constructing an HMM from the family multiple sequence alignment, generating sequences, and extracting the implicit pairwise alignments.

\textbf{Task and Evaluation}: The model is trained to predict the descendant sequence character from the ancestor sequence character given their homologous relationship. We report \textbf{Alignment Accuracy}, defined formally as the per-residue prediction accuracy of the aligned descendant token relative to the ground truth alignment trajectory:
\begin{equation}
    \text{Acc}_{\text{align}} = \frac{1}{|A|} \sum_{(i, j) \in A} \mathbb{I} \Big( \argmax_{c} \hat{Y}_{j,c} = Y_{j} \Big)
\end{equation}
where $A$ is the set of ground-truth aligned residue index pairs $(i, j)$ between the ancestor and descendant, $\hat{Y}_{j,c}$ is the model's predicted probability for amino acid $c$ at descendant position $j$ (which uses the transport plan from ancestor position $i$), and $Y_j$ is the ground-truth amino acid. This directly measures the transport plan's ability to map representations correctly across the evolutionary gap. To isolate the capability of the attention mechanism to learn alignment structure, we chose a deliberately small model scale ($d_{\text{model}}=64, H=4$) that forces the alignment transport plan to perform the heavy lifting.

\textbf{Training}: Models are optimized via Adam with cosine learning rate decay (3e-4 $\to$ 1e-5). To maintain a constant effective batch size across sequence lengths without violating memory constraints, we employ gradient accumulation via \texttt{optax.MultiSteps}. Our implicit differentiation baselines (Section~\ref{sec:implicit_fails}) were thoroughly tuned across learning rates $\in [1e-5, 1e-3]$ and explicit $\alpha$-damping factors $\in [0.1, 0.9]$, but consistently plateaued at $\sim$14\% accuracy.

\textbf{Hardware}: Alignment training was run on Google Cloud TPU v6e-1. Latency micro-benchmarks (Table~\ref{tab:latency}) were collected on TPU v6e-1 (Trillium) with 32\,GB HBM.

\subsection{Results}

% [Table 1 removed for barebones revision]

% [Figure 1 removed for barebones revision]

% [Figure 2 removed for barebones revision]

\subsection{Ablation: Sinkhorn vs.\ Softmax Attention}

As a controlled ablation, we trained an \textbf{identical model} (same architecture, hyperparameters, data, and random seed) replacing only the Sinkhorn attention layer with standard scaled dot-product softmax attention. At 4K tokens, softmax attention plateaus at $\sim$20\% accuracy after 5{,}000 epochs, while Sinkhorn achieves 72\%; a $3.6\times$ improvement (Table~\ref{tab:results}, Figure~\ref{fig:ablation}). The softmax model's predictions degenerate to predicting a single dominant amino acid for all positions, indicating it cannot learn the bijective alignment structure. At 16K tokens, softmax attention exceeds available memory entirely (OOM), while banded Sinkhorn ($W\!=\!2048$) trains to completion.

% [Figure 3 removed for barebones revision]

\subsection{Ablation: Why Multiple Transport Plans Matter}

Table~\ref{tab:results} highlights the critical importance of using distinct transport matrices. All $O(L)$-memory backward methods employing a single transport plan fail to exceed 14\% accuracy, whereas our multi-$P$ approach achieves 72\%.

The failure mode of single-$P$ methods is consistent across runs: training initially improves to $\sim$14\%, then stagnates or collapses entirely with periodic gradient explosions (cross-entropy spikes to $>4.0$). This phenomenon occurs because the single-$P$ approximation introduces a systematic bias in the gradient direction, trapping the optimizer in a poor basin.

\subsection{Scaling to Long Contexts}

% [Table 2 removed for barebones revision]

% [Figure 4 removed for barebones revision]

% [Figure 5 removed for barebones revision]

At 16K tokens, standard softmax attention exceeds available memory entirely, while banded Sinkhorn ($W\!=\!2048$) trains to completion; demonstrating the practical necessity of $O(L{\cdot}W)$ memory scaling for long-context biological alignment (Table~\ref{tab:latency}, Figures~\ref{fig:memory}--\ref{fig:latency}). The banding restriction is biologically well-motivated: homologous residue pairs in protein alignments cluster near the diagonal, making a 2{,}048-residue attention window sufficient to capture relevant correspondences while reducing memory by $8\times$ relative to full attention.

The Pallas kernel enables contexts that are wholly inaccessible to standard methods. Beyond 32K, the naive approach fails entirely (OOM), while banded Pallas Sinkhorn scales cleanly to 131K tokens with linear memory growth.

\section{Related Work}

\textbf{FlashAttention and Tiled Kernels.} \citet{dao2022flashattention} introduced IO-aware tiling for softmax attention. FlashSinkhorn \citep{flashsinkhorn2024} extends this paradigm to Sinkhorn forward passes on GPUs using Triton. Our work differs in three critical ways: (1) we directly address the \textit{backward pass}, which FlashSinkhorn does not tile; (2) we target TPU hardware utilizing Pallas kernels; and (3) we specifically apply Sinkhorn attention to biological sequence alignment rather than point-cloud OT.

\textbf{Differentiable Alignment.} DEDAL \citep{dedal2022} utilizes transformers with differentiable Smith-Waterman for protein alignment. \citet{mensch2018differentiable} proposed smooth relaxations of dynamic programming, while PLASMA \citep{plasma2023} relies on Sinkhorn for protein substructure alignment. Crucially, our method successfully enables alignment at 4$\times$ longer contexts than these approaches.

\textbf{Implicit Differentiation of Fixed Points.} \citet{bai2019deep} demonstrated differentiation through equilibrium models, and \citet{luise2018differential} provided implicit gradients specifically for Sinkhorn. We demonstrate that these approaches fail for Sinkhorn attention due to near-singular adjoint systems, establishing truncated backpropagation as a robust and practical alternative.

\section{Conclusion}

We present the first fully block-wise differentiable Sinkhorn attention mechanism capable of scaling to 16K+ tokens for biological sequence alignment. Our novel multi-transport-plan backward pass resolves a fundamental gap in prior implicit differentiation approaches, achieving 72\% alignment accuracy; a $3.6\times$ improvement over the 20\% achieved by standard softmax attention under identical conditions, and far exceeding the 14\% ceiling of single-$P$ gradient approximations. At 16K tokens, banded Sinkhorn trains to completion where softmax attention exceeds available memory. This approach effectively enables neural sequence alignment at biologically relevant scales that were previously inaccessible to differentiable methods.

\textbf{Limitations and Latency Tradeoffs.} A critical limitation of Sinkhorn attention is its computational latency. As shown in Table~\ref{tab:latency}, our 4K Pallas Sinkhorn forward pass requires $\sim$51 ms compared to 0.26 ms for Flash Softmax—an almost $200\times$ penalty driven by the 15 unrolled Sinkhorn iterations and $O(L)$ block-wise global reductions. The backward pass introduces a further $\sim5\times$ FLOP penalty to recompute the trajectory and apply the distinct transport matrices. Consequently, Sinkhorn attention is best conceptualized as a \textit{high-accuracy offline training tool} for tasks where standard attention fundamentally fails to learn bijective mappings (such as sequence alignment at small embedding dimensions). For real-time inference, a promising future direction is distilling the learned Sinkhorn transport plans into a standard softmax architecture. 

A further consideration is the \textbf{VPU bottleneck}: each Sinkhorn iteration requires elementwise $\exp(\cdot)$ operations, which on modern accelerators execute on multi-function units (MUFU) rather than high-throughput matrix multiply-accumulate (MMA) units. A standard software-hardware optimization established in numerical libraries is \textbf{Cody-Waite range reduction} with polynomial approximation: decomposing $\exp(x) = 2^{k} \cdot p(r)$ where $k = \lfloor x / \ln 2 \rfloor$ is the integer part (computed via bit manipulation) and $p(r)$ is a low-degree polynomial approximation of $\exp(r)$ for $r \in [0, \ln 2)$, executable entirely on FMA units. Implementing this technique in our Sinkhorn kernel would eliminate the MUFU dependency, allowing exponential computation to overlap with MMA operations.

\textbf{Distributed training and banded computation.} Sharding the block-wise Sinkhorn across multiple devices introduces an all-reduce communication step for the potentials $u, v$ after every iteration, since each device holds only a subset of blocks but the logsumexp reduction is global. This per-iteration communication overhead limits scalability. A natural mitigation is \textbf{banded Sinkhorn}: restricting the block-wise computation to a diagonal band of width $W$ around the main diagonal, reducing complexity from $O(L^2)$ to $O(L \cdot W)$. For biological sequences, this is well-motivated since most alignments are approximately diagonal (residue $i$ in the ancestor aligns near residue $i$ in the descendant), and setting $W \ll L$ eliminates cross-device communication for within-band blocks. Crucially, we implement this natively within our custom Pallas backward kernels by dynamically restricting the MXU \texttt{fori\_loop} boundaries to exactly $[i-W, i+W]$ for each query block $i$. This guarantees true $O(L \cdot W)$ hardware utilization and HBM loading without falling back to masking dense operations.

\textbf{Future work.} Future directions include: (1) conducting empirical hardware-level profiling of Cody-Waite polynomial $\exp$ approximations to characterize v6e VPU utilization; (2) replacing position-wise dense layers with 1D convolutions for local motif detection before global alignment; (3) integrating frozen protein language model embeddings (ESM-2) to leverage pre-trained representations; and (4) extending to multiple sequence alignment (MSA) and structural alignment benchmarks.

\bibliographystyle{plainnat}
\bibliography{references}

\appendix

\section{Proof of Proposition~\ref{prop:distinct_P}}
\label{app:proof}

We provide a formal proof in two parts: (A) the four transport matrices are distinct under non-convergence, and (B) collapsing them to a single matrix introduces a quantifiable gradient bias.

\subsection*{Part A: Distinctness of Transport Matrices}

Recall that each transport matrix is defined as:
\begin{equation}
    P^{(a,b)}_{ij} = \exp\!\big(S_{ij} + u^{(a)}_i + v^{(b)}_j\big)
\end{equation}
where $S_{ij}$ is the shared score matrix and $(u^{(a)}, v^{(b)})$ index different potential vectors from the truncated iteration chain.

\begin{lemma}[Distinctness]
Let $\epsilon_u^{(t)} = u^{(t)} - u^*$ and $\epsilon_v^{(t)} = v^{(t)} - v^*$ denote the deviation of the $t$-th iterate from the fixed point $(u^*, v^*)$. If $\|\epsilon_u^{(t)}\|_\infty > 0$ or $\|\epsilon_v^{(t)}\|_\infty > 0$ for any $t \in \{0,1,2\}$, then $P^{(a_1, b_1)} \neq P^{(a_2, b_2)}$ whenever $(a_1, b_1) \neq (a_2, b_2)$.
\end{lemma}

\begin{proof}
Suppose for contradiction that $P^{(a_1, b_1)} = P^{(a_2, b_2)}$ with $(a_1, b_1) \neq (a_2, b_2)$. Then for all $i, j$:
\begin{equation}
    \exp(S_{ij} + u^{(a_1)}_i + v^{(b_1)}_j) = \exp(S_{ij} + u^{(a_2)}_i + v^{(b_2)}_j)
\end{equation}
Taking logarithms:
\begin{equation}
    u^{(a_1)}_i + v^{(b_1)}_j = u^{(a_2)}_i + v^{(b_2)}_j \quad \forall\, i, j
\end{equation}
Fixing $j$ and varying $i$: $u^{(a_1)}_i - u^{(a_2)}_i = c$ for some constant $c$ independent of $i$. Similarly, fixing $i$ and varying $j$: $v^{(b_1)}_j - v^{(b_2)}_j = -c$ for all $j$. Hence $u^{(a_1)} = u^{(a_2)} + c\mathbf{1}$ and $v^{(b_1)} = v^{(b_2)} - c\mathbf{1}$.

Now consider $a_1 \neq a_2$. Without loss of generality, assume $a_1 > a_2$, which implies $u^{(a_1)}$ occurs later in the iteration chain than $u^{(a_2)}$. This requires $u^{(a_1)} = u^{(a_2)} + c\mathbf{1}$. Substituting into the Sinkhorn update rules for both iterates:
\begin{equation}
    -\log\sum_j \exp(S_{ij} + v^{(b_1)}_j) = -\log\sum_j \exp(S_{ij} + v^{(b_2)}_j) + c
\end{equation}
This implies $\sum_j \exp(S_{ij} + v^{(b_1)}_j) = e^{-c} \sum_j \exp(S_{ij} + v^{(b_2)}_j)$ for all $i$. For this to hold across arbitrary score matrices $S$, it requires $v^{(b_1)} = v^{(b_2)} - c\mathbf{1}$, which means the $v$ potential has fundamentally converged up to a global uniform shift. 

Crucially, while standard doubly-stochastic Sinkhorn is shift-invariant to $(u+c\mathbf{1}, v-c\mathbf{1})$, the \textit{unbalanced} Sinkhorn with a fixed gap score $g$ breaks this shift invariance. The potential update includes the constant $\exp(g)$:
\begin{equation}
    \sum_j \exp(S_{ij} + v^{(b_1)}_j) + \exp(g) = e^{-c} \Big( \sum_j \exp(S_{ij} + v^{(b_2)}_j) + \exp(g+c) \Big)
\end{equation}
For this equation to hold, the gap term requires $e^g = e^{g+c}$, forcing $c=0$. Therefore, if $P^{(a_1, b_1)} = P^{(a_2, b_2)}$, then $u^{(a_1)} = u^{(a_2)}$ and $v^{(b_1)} = v^{(b_2)}$, contradicting the premise that the potentials have not converged.
\end{proof}

\subsection*{Part B: Gradient Bias Under Single-$P$ Assumption}

Let $\nabla_S \mathcal{L}_{\text{multi}}$ denote the true multi-$P$ gradient and $\nabla_S \mathcal{L}_{\text{single}}$ the approximation using a single $P = P^{(2,2)}$ for all steps. Define the bias $\mathcal{B} = \nabla_S \mathcal{L}_{\text{single}} - \nabla_S \mathcal{L}_{\text{multi}}$.

\begin{theorem}[Gradient Bias Bound]
Let $\Delta u^{(t)} = u^{(t+1)} - u^{(t)}$ and $\Delta v^{(t)} = v^{(t+1)} - v^{(t)}$ be the per-iteration potential updates. The gradient bias satisfies:
\begin{equation}
    \mathcal{B}_{mn} = \sum_{k \in \{C,D,E\}} \Big(P^{(2,2)}_{mn} - P^{(a_k, b_k)}_{mn}\Big) \cdot \gamma^{(k)}_m
\end{equation}
where $\gamma^{(k)}$ are the corresponding gradient scalars ($\delta u^{(2)}, \delta v^{(1)}, \delta u^{(1)}$ respectively), and:
\begin{equation}
    P^{(2,2)}_{mn} - P^{(a_k, b_k)}_{mn} = P^{(2,2)}_{mn}\Big(1 - \exp\!\big((u^{(a_k)} - u^{(2)})_m + (v^{(b_k)} - v^{(2)})_n\big)\Big)
\end{equation}
\end{theorem}

\begin{proof}
The total gradient is $\nabla_S \mathcal{L} = \sum_{k \in \{A,\ldots,E\}} \Delta S^{(k)}$, where each term has the form $\Delta S^{(k)}_{mn} = \gamma^{(k)} \cdot P^{(a_k, b_k)}_{mn}$ (with appropriate sign). The single-$P$ approximation replaces each $P^{(a_k, b_k)}$ with $P^{(2,2)}$, leaving terms $A$ and $B$ unchanged (they already use $P^{(2,2)}$) but biasing terms $C$, $D$, and $E$.

For term $C$: $\Delta S^{(C)}_{mn} = -\delta u^{(2)}_m \cdot P^{(2,1)}_{mn}$. The single-$P$ version uses $P^{(2,2)}_{mn}$ instead. Since:
\begin{equation}
    \frac{P^{(2,2)}_{mn}}{P^{(2,1)}_{mn}} = \exp\big(v^{(2)}_n - v^{(1)}_n\big) = \exp(\Delta v^{(1)}_n)
\end{equation}
the bias contribution from term $C$ is:
\begin{equation}
    \mathcal{B}^{(C)}_{mn} = -\delta u^{(2)}_m \cdot P^{(2,1)}_{mn} \cdot \big(\exp(\Delta v^{(1)}_n) - 1\big)
\end{equation}

Similarly, for term $E$ (the deepest in the chain):
\begin{equation}
    \frac{P^{(2,2)}_{mn}}{P^{(1,0)}_{mn}} = \exp\big(\Delta u^{(1)}_m + \Delta v^{(0)}_n + \Delta v^{(1)}_n\big)
\end{equation}
which accumulates errors from \textit{all} intermediate potential updates.

Critically, $\|\mathcal{B}\|$ does not vanish because in practice the forward Sinkhorn iterations are truncated at a finite $T=15$. The potential gaps $\Delta u^{(t)}, \Delta v^{(t)}$ thus remain non-zero for the backward refinement iterations, making the $\exp(\Delta v^{(1)}_n) - 1$ factor significant. This non-vanishing bias explains the empirical observation that single-$P$ methods inevitably plateau at 14\% accuracy.
\end{proof}

\begin{remark}
The bias is \textit{multiplicative}, not additive: it scales with $P^{(a_k,b_k)}_{mn}$ itself. Large transport entries (high-confidence alignments) receive the largest bias, systematically corrupting the gradient signal precisely where it matters most for learning.
\end{remark}

\end{document}